\documentclass[12pt,a4paper]{article}

\usepackage{template}
\usepackage{amsmath}
\usepackage{amssymb}
\usepackage{booktabs}

\hypersetup{
  colorlinks=true,
  linkcolor=blue!45!black,
  urlcolor=blue!45!black,
  citecolor=blue!45!black,
  pdfauthor={Иван Афанасьев},
  pdftitle={Задача 2. Перекрывающиеся прямоугольники (элементарная версия)}
}

\begin{document}
\sloppy

\begin{center}
  {\LARGE\bfseries Задача 2. Перекрывающиеся прямоугольники}\\[4pt]
  {\large элементарная версия (только школьные средства)}\\[6pt]
  {\large Иван Афанасьев}\\[2pt]
  {\normalsize Вступительный набор <<group2026famcs>>, ФПМИ БГУ, 2026}
\end{center}

\vspace{6pt}
\begin{center}
\fbox{\begin{minipage}{0.92\textwidth}
\small
\textbf{Использование ИИ.} Решение подготовлено при активном участии
\textbf{GLM-5.2} (помощь в проверке и оформлении гипотез) и
\textbf{OpenCode-CI-Agent} Выполнял независимую рецензию на изменения во время Pull Request'а: нашёл ошибку в подсчёте для обобщения~4 (в счётчике последовательного
режима была индексная ошибка; неверное число $25$ заменено верным $16$).
\end{minipage}}
\end{center}
\vspace{10pt}

\section{Постановка задачи}

На поле $n\times n$ по линиям клеток выкладывают одну за другой цветные
костяшки $1\times2$ (можно поворачивать на $90^{\circ}$, выходить за поле
нельзя, класть поверх уже лежащих~--- можно). В клетке виден цвет
\emph{последней} накрывшей её костяшки; клетка, которую не накрыла ни одна
костяшка, остаётся пустой (фоном). Дана итоговая картина. Режимы:
\begin{enumerate}[label=\arabic*.]
  \item дан \emph{запас} костяшек ($a_1$ штук цвета $1$, $a_2$ цвета $2$,
        \dots); выложить нужно все, порядок любой. Найти порядок ходов,
        дающий картину, или обосновать, что его нет;
  \item дана \emph{последовательность} цветов $c_1,\dots,c_m$~--- костяшки
        обязаны выкладываться именно в этом порядке (позиции выбираем
        сами); та же задача.
\end{enumerate}
Обобщения 3--4: посчитать, сколько различных итоговых картин вообще можно
получить из данного запаса (или последовательности). По уточнению авторов:
$m\le100$; в наших примерах используется $n\le10$, $k\le9$, $a_i\le10$.

Все режимы реализованы программой \texttt{overlap.cpp} (ZIP-архив);
программа печатает последовательность ходов и сама проверяет её,
моделируя выкладывание заново.

\section{Главная идея: снимаем костяшки с конца}

Смотреть на игру удобно \emph{задом наперёд}: будем не класть костяшки, а
снимать их, начиная с самой последней.

\begin{itemize}
  \item Самая последняя костяшка лежит поверх всех, поэтому \emph{обе} её
        клетки на картине показывают её цвет.
  \item Когда мы её снимаем, про две освободившиеся клетки мы уже знаем,
        что их цвет <<объяснён>>: что бы под ними ни оказалось, позже это
        всё равно закрасится. Назовём такие клетки \emph{свободными}.
  \item Клетки фона нельзя накрывать вообще ни разу: любая накрывшая их
        костяшка осталась бы видна там, где должен быть фон.
\end{itemize}

Итого правило снятия: костяшку цвета $c$ можно снять с пары соседних
нефоновых клеток, если каждая из этих двух клеток либо уже свободна, либо
показывает цвет $c$. После снятия обе клетки свободны. Картина получается
в игре тогда и только тогда, когда можно снять весь набор костяшек так,
чтобы каждая нефоновая клетка хотя бы раз была накрыта (стала свободной).

\paragraph{Порядок снятия не важен.} Ключевое наблюдение: снятие любой
костяшки лишь превращает клетки в свободные, а от этого другие снятия
могут только \emph{стать возможными}, но никогда не становятся
невозможными. Отсюда два следствия.
\begin{itemize}
  \item \emph{Проверка достижимости.} Будем жадно снимать что снимается
        (цветов и числа костяшек пока не ограничиваем). Если поле
        сняли целиком~--- картина достижима; если жадный процесс
        застрял~--- недостижима, в каком бы порядке мы ни действовали
        (застрять <<по своей вине>> нельзя: всё, что было снимаемо,
        осталось снимаемым).
  \item \emph{Свобода выбора.} Если картина достижима, то любой
        разумный порядок снятия можно доводить до конца~--- это
        обосновывает работу перебора ниже.
\end{itemize}
Например, у достижимой непустой картины обязательно есть две соседние
клетки одного цвета (след самой последней костяшки)~--- удобный экспресс-тест.

\section{Режим 1: запас костяшек}

\paragraph{Сколько костяшек цвета $c$ нужно обязательно.} Пусть $S_c$~---
клетки цвета $c$ на картине. Разобьём как можно больше клеток из $S_c$ на
пары \emph{соседних} клеток (пары не пересекаются); пусть максимально
возможное число таких пар равно $\nu_c$. Тогда минимально необходимое
число костяшек цвета $c$ равно
\[
  \mathrm{need}_c=|S_c|-\nu_c .
\]
\emph{Почему не меньше.} Каждая клетка из $S_c$ обязана быть накрыта в
последний раз костяшкой цвета $c$. Одна костяшка может быть <<последней>>
не более чем для двух клеток, и если ровно для двух~--- то это две
соседние клетки цвета $c$, т.\,е. одна из наших пар. Значит, если костяшек
цвета $c$ меньше, чем $|S_c|-\nu_c$, все клетки цвета $c$ показать
невозможно.

\emph{Почему достаточно.} Пары снимаем первыми (в игре они кладутся
последними): каждая пара~--- одна костяшка, обе клетки показывают её
цвет, снятие законно. Оставшиеся <<одиночки>> закрываем по одной костяшке:
одиночку $u$ цвета $c$ снимаем вместе с каким-нибудь нефоновым соседом,
который к этому моменту уже свободен. Застрять здесь нельзя: если бы в
какой-то момент ни одну одиночку нельзя было снять, то у каждой
оставшейся одиночки все нефоновые соседи~--- тоже оставшиеся одиночки.
Возьмём ту из них, которая при жадной проверке достижимости снимается
первой; её партнёр по снятию~--- сосед-одиночка, ещё не снятый, значит
снятие возможно лишь при совпадении цветов. Но две соседние одиночки
одного цвета можно было бы объединить в пару~--- вопреки максимальности
числа пар $\nu_c$. Противоречие.

\paragraph{Критерий и алгоритм.} Итого режим~1 разрешим тогда и только
тогда, когда (а)~картина достижима (жадная проверка), (б)~$a_c\ge
\mathrm{need}_c$ для каждого цвета (для пустой картины~--- когда $m=0$).
Излишек $a_c-\mathrm{need}_c$ прячем: эти костяшки кладём самыми первыми
на любую пару соседних нефоновых клеток~--- их потом закрасят. Программа
строит порядок именно так: снимает костяшки с конца, каждый раз выбирая
самую <<стеснённую>> клетку (у которой меньше всего вариантов снятия);
максимальное число пар $\nu_c$ она находит стандартным перебором с
улучшениями (метод чередующихся цепочек; клетки поля можно раскрасить в
шахматном порядке, что ускоряет поиск). В конце программа проигрывает
найденные ходы заново и сверяет результат с картиной
(строка \texttt{Self-check ... PASS}).

\section{Режим 2: заданная последовательность цветов}

Снимаем строго в обратном порядке: сначала костяшку цвета $c_m$, затем
$c_{m-1}$ и т.\,д. На каждом шаге программа перебирает все допустимые
позиции (обе клетки свободны или показывают нужный цвет); некоторые
костяшки приходится <<прятать>> на уже свободные клетки~--- перебор
учитывает и это. Чтобы перебор не разрастался, используются простые
отсечения:
\begin{itemize}
  \item клеток цвета $c$, ещё не объяснённых, не может быть больше, чем
        $2\times$(число оставшихся костяшек цвета $c$);
  \item не должно оставаться <<замурованных>> клеток без нефоновых
        соседей;
  \item состояния, из которых решение уже не нашлось, запоминаются и
        повторно не разбираются.
\end{itemize}
Перебор полный, поэтому ответ <<невозможно>> достоверен. Найденный
порядок также проверяется симуляцией.

\section{Обобщения 3--4: сколько бывает итоговых картин}

Считаем прямым перебором: кладём костяшки по одной всеми возможными
способами (в режиме запаса~--- любой оставшийся цвет, в режиме
последовательности~--- строго очередной) и складываем все получившиеся
картины в множество без повторов. Для поля $2\times2$ и запаса
$\{1{:}1,\,2{:}1\}$ получается $28$ различных картин, а для той же пары
костяшек в фиксированном порядке <<сначала~1, потом~2>>~--- $16$:
остальные $12$ картин требуют положить костяшку $2$ раньше костяшки $1$
(это картины, где цвет $1$ виден <<поверх>> цвета $2$ либо костяшка
цвета~$2$ скрыта целиком). Оба числа сверены с независимым полным
перебором всех партий. Для больших полей числа растут стремительно,
поэтому в программе предусмотрен ограничитель размера перебора с честным
сообщением вместо зависания.

\section{Проверка корректности}

\begin{itemize}
  \item \textbf{Примеры.} Точный запас; запас с излишком (лишняя костяшка
        прячется); недостаточный запас (программа называет цвет и сколько
        именно не хватает); правильный и неправильный порядок в режиме~2.
  \item \textbf{Рисунок из условия.} Итоговая конфигурация с~$14$
        костяшками из условия записана в \texttt{examples/pdf\_pattern.txt};
        программа восстанавливает её из минимального набора ($8$ костяшек),
        самопроверка проходит.
  \item \textbf{Случайные испытания.} Скрипт \texttt{stress.py} играет
        $500$ случайных партий (значит, их картины заведомо достижимы) и
        проверяет, что оба режима находят решение; ещё $200$ произвольных
        картин проверяются на согласованность ответа с самопроверкой.
        Сбоев нет.
  \item \textbf{Полная сверка.} Дополнительно ответы обоих режимов были
        сверены с независимым полным перебором \emph{всех} партий на полях
        $2\times2$ и $3\times3$ для нескольких запасов и
        последовательностей~--- совпадение на всех тысячах картин, в обе
        стороны (и <<возможно>>, и <<невозможно>>).
\end{itemize}

\section{Приложение: сборка и запуск}

\begin{verbatim}
g++ -O2 -std=c++17 -o overlap overlap.cpp

# Режим 1: запас костяшек a1..ak; картина подаётся на вход
./overlap mode1 7 7 2 0 1 1 2 1 1 < examples/pdf_pattern.txt

# Режим 2: фиксированная последовательность из m цветов
printf "1 2 2\n0 0 0\n0 0 0\n" | ./overlap mode2 3 2 1 2

# Обобщения 3-4: число различных итоговых картин (малые n)
./overlap count1 2 2 1 1
./overlap count2 2 2 1 2
\end{verbatim}
Формат картины: $n$ строк по $n$ чисел, $0$~--- фон, $1,\dots,k$~---
цвета. Вывод~--- ходы в порядке выкладывания (первые ходы~--- нижние
костяшки) и строка самопроверки.

\end{document}
